% Chapter 3, Figure 44: the gross planform area sigma_F of six fin shapes (scan: figures/ch3/fig44.png; shapes
% after G. Harry Stine). The 1973 artwork is a ruled 2 x 3 grid of named cells (a table drawn in the artwork,
% STYLE.md section 14): set here as a table with booktabs rules (top, between the rows, bottom; no vertical rules
% or boxes, as section 8), the names under the drawings. "Aerobee-Hi" as the text and caption spell it: the
% hand's last glyph is a dotless cap-height stroke, which is how it writes i (the i of "Trapezoidal" here), so the
% print does not show the case (as the Figure 41 redraw).
% Each cell: half of the rocket seen from the side, nose up (as Figure 49): the body tube broken off at the top
% (\breakline), its centre line, one fin on the right; the gross area (the exposed fin and its extension into the
% tube to the centre line) hatched, the leading and trailing edges extended to the centre line dashed (guide). For
% the Python-2 the lines are drawn from the root chord's ends parallel to the base (the text, ch3-sec6a.tex), the
% trailing one along the base itself, and the fin, which runs aft past the base, is broken off.
% Coordinates are the scan's pixels (150 per inch) from the centre line at the base, y up, drawn 1.3 times the
% printed size; one tube for all six (radius 18.5, 101 long), the fins as measured: every outline lies within
% 3 px of the scan's ink (95th percentiles 0.2-2.0 px; Swept and Python-2 made exactly parallel-edged).
\documentclass[11pt]{standalone}
\usepackage{tamrfig}
\begin{document}
\begin{tikzpicture}[x=0.022cm, y=0.022cm, break amplitude=0.36]
  \def\R{18.5}\def\T{101}            % tube radius and length
  \def\W{180}\def\H{194}             % cell width and height
  % the tube, its break and centre line (drawn over the hatching)
  \newcommand{\tube}{%
    \draw[outline] (-\R,\T) -- (-\R,0) -- (\R,0) -- (\R,\T);
    \breakline{(-\R,\T)}{(\R,\T)}
    \draw[centerline] (0,-14) -- (0,\T+10);}
  % the same zigzag as \breakline (at the break amplitude 0.36 set above), as a path to bound a fill:
  % \breakpath{<from>}{<to>}
  \newcommand{\breakpath}[2]{%
    let \p1=($#2-#1$), \n1={veclen(\x1,\y1)}, \n2={atan2(\y1,\x1)} in
    \foreach \u/\v in {0.172/0, 0.276/0.241, 0.414/-0.483, 0.552/0.5, 0.664/-0.241, 0.776/0}
      { -- ($#1+(\n2:{\u*\n1*1pt})+(\n2+90:{\v*0.36*\n1*1pt})$) } -- #2}
  % the name under a cell's drawing, on one baseline
  \newcommand{\cellname}[1]{\node[anchor=base, font=\small] at (46,-55) {#1};}
  % ---- row 1 --------------------------------------------------------------------------------------------
  \begin{scope}[shift={(0.5*\W-46,0)}]   % Delta
    \fill[hatch] (0,84.55) -- (\R,70) -- (107.5,0) -- (0,0) -- cycle;
    \tube
    \draw[outline] (\R,70) -- (107.5,0) -- (\R,0);
    \draw[guide] (\R,70) -- (0,84.55);
    \cellname{Delta}
  \end{scope}
  \begin{scope}[shift={(1.5*\W-46,0)}]   % Rectangular
    \fill[hatch] (0,70.5) rectangle (103.5,0);
    \tube
    \draw[outline] (\R,70.5) -- (103.5,70.5) -- (103.5,0) -- (\R,0);
    \draw[guide] (\R,70.5) -- (0,70.5);
    \cellname{Rectangular}
  \end{scope}
  \begin{scope}[shift={(2.5*\W-46,0)}]   % Swept: leading and trailing edges parallel
    \fill[hatch] (0,79.04) -- (\R,71) -- (110.5,31) -- (110.5,-40) -- (\R,0) -- (0,8.04) -- cycle;
    \tube
    \draw[outline] (\R,71) -- (110.5,31) -- (110.5,-40) -- (\R,0);
    \draw[guide] (\R,71) -- (0,79.04) (\R,0) -- (0,8.04);
    \cellname{Swept}
  \end{scope}
  % ---- row 2 --------------------------------------------------------------------------------------------
  \begin{scope}[shift={(0.5*\W-46,-\H)}]   % Trapezoidal: the trailing edge extended meets the centre line below the base
    \fill[hatch] (0,79.04) -- (\R,75) -- (105.5,56) -- (105.5,18) -- (\R,0) -- (0,-3.83) -- cycle;
    \tube
    \draw[outline] (\R,75) -- (105.5,56) -- (105.5,18) -- (\R,0);
    \draw[guide] (\R,75) -- (0,79.04) (\R,0) -- (0,-3.83);
    \cellname{Trapezoidal}
  \end{scope}
  \begin{scope}[shift={(1.5*\W-46,-\H)}]   % Aerobee-Hi: the 45-degree chamfer at the root extended
    \fill[hatch] (0,85.59) -- (\R,75) -- (101.5,27.5) -- (101.5,-18.5) -- (37,-18.5) -- (\R,0) -- (0,\R) -- cycle;
    \tube
    \draw[outline] (\R,75) -- (101.5,27.5) -- (101.5,-18.5) -- (37,-18.5) -- (\R,0);
    \draw[guide] (\R,75) -- (0,85.59) (\R,0) -- (0,\R);
    \cellname{Aerobee-Hi}
  \end{scope}
  \begin{scope}[shift={(2.5*\W-46,-\H)}]   % Python-2: edges parallel, broken off 26 below the base
    \fill[hatch] (0,0) -- (\R,0) -- (33.6,-26) \breakpath{(33.6,-26)}{(77.3,-26)} -- (25.5,63) -- (0,63) -- cycle;
    \tube
    \draw[outline] (\R,63) -- (25.5,63) -- (77.3,-26);
    \draw[outline] (\R,0) -- (33.6,-26);
    \breakline{(33.6,-26)}{(77.3,-26)}
    \draw[guide] (0,63) -- (\R,63);
    \cellname{Python-2}
  \end{scope}
  % ---- rules (booktabs: \heavyrulewidth 0.08em, \lightrulewidth 0.05em) --------------------------------------
  \draw[ink, line width=0.876pt] (0,123) -- (3*\W,123);
  \draw[ink, line width=0.548pt] (0,123-\H) -- (3*\W,123-\H);
  \draw[ink, line width=0.876pt] (0,123-2*\H) -- (3*\W,123-2*\H);
\end{tikzpicture}
\end{document}
